\scanpage{81}
\[
\left\{\begin{aligned}
&\text{paires de pantalons \deleted{\unsure{dég|dégé}} généralisés \inserted{$X',X''$ à \unsure{trous|bord}}, indexés}\\
&\text{par }E'\amalg\{s\},\ E''\amalg\{t\},\ \text{avec }s,t\text{ correspondant}\\
&\text{aux \deleted{bords} composantes ponctuelles du}\\
&\text{bord, avec les paramètres }(\lambda_i)_{i\in E'\amalg E''}\in\mathbb R^{\unsure{*+E'\amalg E''|\Sigma'\amalg\Sigma''}},\\
&\text{et un isom d'identification}\\
&\text{dans }\mathbb C^*\wedge_{\pm1}(E'\wedge E'')
\end{aligned}\right.
\]

Une fois ce cas compris, et revenant alors au cas général, il faudrait déjà, pour un type de découpage donné incarné par un diagramme $\mathrm{MD}$ du $md_{g,\nu}$, comprendre l'application sur les groupes fondamentaux induite par
\[
\mathrm{MDB}_{\unsure{g}}\longrightarrow \mathrm M_{\mathrm{st}}\mathrm{ND}_{g,\nu}, \tag*{($*$)}
\]
celui des deuxièmes membres étant essentiellement compris en termes de groupes de Teichmüller spéciaux -- \uline{mais il doit en être pareil pour les premiers} ! Également, il ne devrait pas être bien dur de vérifier si ($*$)
% archive sheet number: 83
